\section{Mentoring \& Inclusion}
An outstanding educator must also be a supportive mentor who actively fosters an inclusive environment. In my academic career, I strive to dismantle systemic barriers in STEM through intentional classroom design and research mentorship:

\subsection{Creating an Inclusive Classroom Environment}
To ensure all students feel valued, I employ inclusive teaching practices, such as using gender-neutral and culturally diverse examples in assignments. I implement anonymous Q\&A boards (e.g., Ed Discussion) so students can ask questions without fear of judgment. I also use a transparent ``slip day'' policy for assignments, which accommodates personal, health, or financial challenges without requiring students to justify extensions.

\subsection{Mentoring Undergraduate Researchers}
Over the past four years, I have mentored 5 undergraduate students through Stanford's CURIS research program, focusing on underrepresented backgrounds. I meet with mentees weekly, providing feedback on research and professional development. Two have co-authored papers with me at CVPR and ICML, and three have pursued computer science PhDs. One former mentee noted: \emph{``Alex did not just hand me tasks; he taught me how to read papers critically, write code, and present with confidence.''}

% --- Section: Future Teaching Goals ---
